% ==============================================================================
% CHAPITRE 8 : STATISTIQUES À UNE VARIABLE (VERSION ÉLÈVE)
% ==============================================================================
\chapter{Statistiques à une Variable}

% ------------------------------------------------------------------------------
% 1. ACTIVITÉ D'INVESTIGATION (SITUATION PROFESSIONNELLE : INDUSTRIE / CONTRÔLE QUALITÉ)
% ------------------------------------------------------------------------------
\begin{activite}{Contrôle Qualité sur Pièces Usinées (Industrie \& Mécanique)}
Un technicien en contrôle qualité mesure au pied à coulisse le diamètre (en $\text{mm}$) d'un échantillon de \textbf{20 pièces} métalliques usinées en série :
\begin{center}
\renewcommand{\arraystretch}{1.2}
\begin{tabular}{|c|c|c|c|c|c|c|c|c|c|}
\hline
49{,}8 & 50{,}1 & 49{,}9 & 50{,}0 & 50{,}2 & 49{,}9 & 50{,}0 & 50{,}1 & 49{,}8 & 50{,}0 \\
\hline
50{,}3 & 49{,}9 & 50{,}1 & 50{,}0 & 49{,}8 & 50{,}2 & 50{,}0 & 50{,}1 & 49{,}9 & 50{,}0 \\
\hline
\end{tabular}
\end{center}

\vspace{1mm}
\noindent
\begin{minipage}[t]{0.53\linewidth}
\textbf{Tableau des effectifs regroupés :}
\begin{center}
\renewcommand{\arraystretch}{1.2}
\setlength{\tabcolsep}{2pt}
\begin{tabularx}{\linewidth}{|l|*{6}{>{\centering\arraybackslash}X|}}
\hline
\rowcolor{slate800} \textcolor{white}{\textbf{Diamètre}} & \textcolor{white}{\scriptsize 49{,}8} & \textcolor{white}{\scriptsize 49{,}9} & \textcolor{white}{\scriptsize 50{,}0} & \textcolor{white}{\scriptsize 50{,}1} & \textcolor{white}{\scriptsize 50{,}2} & \textcolor{white}{\scriptsize 50{,}3} \\
\hline
\textbf{Effectif} & 3 & \pointille[0.5cm] & 6 & \pointille[0.5cm] & 2 & 1 \\
\hline
\end{tabularx}
\end{center}
\end{minipage}\hfill
\begin{minipage}[t]{0.44\linewidth}
\begin{tcolorbox}[colback=slate100, colframe=colPrimary, arc=4pt, boxrule=0.8pt, top=4pt, bottom=4pt, left=6pt, right=6pt]
\sffamily\footnotesize
\textbf{Problématique :}
\begin{enumerate}[leftmargin=1.2em, itemsep=2pt]
    \item Quel est le diamètre le plus fréquent ?
    \item Quelle est la moyenne de cet échantillon ?
    \item Comment savoir si la production est conforme au cahier des charges ($50{,}00\text{ mm} \pm 0{,}25\text{ mm}$) ?
\end{enumerate}
\end{tcolorbox}
\end{minipage}

\vspace{2mm}
\noindent\textbf{Espace de recherche :}
\lignesreponse[3]
\end{activite}

% ------------------------------------------------------------------------------
% 2. COURS : VOCABULAIRE, EFFECTIFS & FRÉQUENCES
% ------------------------------------------------------------------------------
\section{Vocabulaire, Effectifs \& Fréquences}

\begin{cadredef}[Vocabulaire Statistique]
\begin{itemize}[leftmargin=1.5em, itemsep=2pt]
    \item \textbf{Population} : L'ensemble des individus étudiés (ex : l'ensemble des pièces produites).
    \item \textbf{Caractère} : La grandeur mesurée ou observée (ex : le diamètre en mm, la couleur, le salaire).
    \item \textbf{Effectif total ($N$)} : Le nombre total d'individus dans l'échantillon ($N = \sum n_i$).
    \item \textbf{Fréquence ($f$)} : La proportion d'un effectif par rapport à l'effectif total :
    \[ \mathbf{f = \frac{\text{Effectif } n_i}{\text{Effectif total } N}} \qquad \text{et en pourcentage : } \mathbf{f\% = \frac{n_i}{N} \times 100} \]
\end{itemize}
\end{cadredef}

\begin{cadreexemple}
Dans un groupe de $25$ élèves, $10$ sont en 3PM A.
\begin{itemize}[leftmargin=1.5em, itemsep=2pt]
    \item Fréquence : $f = \frac{10}{25} = \pointille[1.5cm]$
    \item Fréquence en pourcentage : $f\% = \frac{10}{25} \times 100 = \pointille[1.5cm] \text{ \%}$
\end{itemize}
\end{cadreexemple}

% ------------------------------------------------------------------------------
% 3. COURS : INDICATEURS DE POSITION (MOYENNE & MÉDIANE)
% ------------------------------------------------------------------------------
\section{Indicateurs de Position : Moyenne \& Médiane}

\begin{cadremethode}[Calculer une Moyenne Simple et une Moyenne Pondérée]
\begin{enumerate}[leftmargin=1.5em, itemsep=3pt]
    \item \textbf{Moyenne simple $\bar{x}$} : On additionne toutes les valeurs et on divise par le nombre total de valeurs :
    \[ \bar{x} = \frac{x_1 + x_2 + \dots + x_p}{N} \]
    \item \textbf{Moyenne pondérée $\bar{x}$} : On multiplie chaque valeur par son effectif, on fait la somme, puis on divise par l'effectif total $N$ :
    \[ \mathbf{\bar{x} = \frac{n_1 \times x_1 + n_2 \times x_2 + \dots + n_p \times x_p}{N}} \]
\end{enumerate}
\end{cadremethode}

\begin{cadredef}[Médiane d'une Série Statistique]
La \textbf{médiane} (notée $Me$) d'une série ordonnée de valeurs est le nombre qui partage cette série en \textbf{deux groupes de même effectif} :
\begin{itemize}[leftmargin=1.5em, itemsep=2pt]
    \item Au moins $50\text{ \%}$ des valeurs sont inférieures ou égales à la médiane.
    \item Au moins $50\text{ \%}$ des valeurs sont supérieures ou égales à la médiane.
\end{itemize}
\end{cadredef}

\begin{cadremethode}[Méthode pas à pas : Déterminer la Médiane]
\textbf{Étape 1} : On ordonne impérativement les valeurs dans l'\textbf{ordre croissant}.\\
\textbf{Étape 2} : On calcule le rang médian selon la parité de l'effectif total $N$ :
\begin{itemize}[leftmargin=1.5em, itemsep=2pt]
    \item \textbf{Si $N$ est impair} (ex : $N = 9$) : $\frac{N+1}{2} = \frac{10}{2} = 5 \implies$ La médiane est la \textbf{5\textsuperscript{ème} valeur}.
    \item \textbf{Si $N$ est pair} (ex : $N = 10$) : $\frac{N}{2} = 5 \implies$ La médiane est la \textbf{moyenne entre la 5\textsuperscript{ème} et la 6\textsuperscript{ème} valeur}.
\end{itemize}
\end{cadremethode}

% ------------------------------------------------------------------------------
% 4. COURS : INDICATEUR DE DISPERSION (ÉTENDUE)
% ------------------------------------------------------------------------------
\section{Indicateur de Dispersion : L'Étendue}

\begin{cadredef}[Étendue ($e$)]
L'\textbf{étendue} d'une série statistique est la différence entre la plus grande valeur ($x_{\text{max}}$) et la plus petite valeur ($x_{\text{min}}$) :
\[ \mathbf{e = x_{\text{max}} - x_{\text{min}}} \]
\textit{Plus l'étendue est petite, plus la série est homogène et régulière.}
\end{cadredef}

\begin{cadreretenu}[L'Essentiel du Chapitre 8]
\begin{itemize}[leftmargin=1.5em, itemsep=2pt]
    \item \textbf{Fréquence} : $f = \frac{n_i}{N}$ (la somme des fréquences vaut toujours $1$ ou $100\%$).
    \item \textbf{Moyenne $\bar{x}$} : somme de (valeur $\times$ effectif) divisée par l'effectif total.
    \item \textbf{Médiane $Me$} : valeur centrale de la série triée (50\% au-dessus, 50\% en dessous).
    \item \textbf{Étendue $e$} : $x_{\text{max}} - x_{\text{min}}$ (mesure l'écartement des valeurs).
\end{itemize}
\end{cadreretenu}

% ------------------------------------------------------------------------------
% 5. QUESTIONS FLASH / AUTOMATISMES
% ------------------------------------------------------------------------------
\section{Questions Flash / Automatismes}

\begin{automatisme}[8 Questions Flash d'Entraînement]
\begin{enumerate}[label=\textbf{Q\arabic*.}]
    \item Calculer la moyenne des notes : $12$ ; $14$ ; $16$ $\implies \bar{x} = \pointille[3cm]$
    \item Donner l'étendue de la série : $4$ ; $9$ ; $15$ ; $21$ $\implies e = \pointille[3cm]$
    \item Trouver la médiane de la série ordonnée : $3$ ; $7$ ; \textbf{11} ; $14$ ; $19$ $\implies Me = \pointille[3cm]$
    \item Trouver la médiane de la série ordonnée : $6$ ; $8$ ; $12$ ; $16$ $\implies Me = \pointille[3cm]$
    \item Dans un atelier de 20 ouvriers, 5 sont apprentis. Fréquence en $\%$ : \pointille[3cm]
    \item Si l'effectif total est $N = 50$ et l'effectif d'une modalité est $15$, la fréquence vaut : \pointille[3cm]
    \item Une série a pour minimum $12$ et pour étendue $18$. Quelle est la valeur maximale ? \pointille[3cm]
    \item Vrai ou Faux : La médiane est toujours égale à la moyenne. \pointille[3cm]
\end{enumerate}
\end{automatisme}

% ------------------------------------------------------------------------------
% 6. TRAVAUX DIRIGÉS (TD) - EXERCICES D'ENTRAÎNEMENT
% ------------------------------------------------------------------------------
\section{Travaux Dirigés (TD) : Exercices d'Entraînement}

\begin{exercice}[2 pts]{\capRealiser}{Série brute et indicateurs de base}
Voici les temps d'attente (en minutes) relevés pour 7 clients au service après-vente :
\[ 8 \quad ; \quad 14 \quad ; \quad 5 \quad ; \quad 12 \quad ; \quad 20 \quad ; \quad 6 \quad ; \quad 12 \]
\begin{enumerate}[leftmargin=1.5em]
    \item Rallonger la série dans l'ordre croissant : \pointille[7cm]
    \item Calculer le temps moyen d'attente $\bar{x}$.
    \item Déterminer le temps médian d'attente $Me$.
    \item Calculer l'étendue $e$ de cette série.
\end{enumerate}
\lignesreponse[3]
\end{exercice}

\begin{exercice}[3 pts]{\capRealiser}{Tableau d'effectifs \& Moyenne pondérée (Restauration)}
Un restaurant traditionnel note le nombre de couverts servis chaque midi pendant un mois (25 jours d'ouverture) :
\begin{center}
\renewcommand{\arraystretch}{1.2}
\begin{tabularx}{0.9\linewidth}{|l|*{5}{>{\centering\arraybackslash}X|}}
\hline
\rowcolor{slate800} \textcolor{white}{\textbf{Nombre de couverts ($x_i$)}} & \textcolor{white}{30} & \textcolor{white}{40} & \textcolor{white}{50} & \textcolor{white}{60} & \textcolor{white}{70} \\
\hline
\textbf{Nombre de jours ($n_i$)} & 3 & 7 & 9 & 4 & 2 \\
\hline
\textbf{Fréquence en \%} & \pointille[1cm] & \pointille[1cm] & \pointille[1cm] & \pointille[1cm] & \pointille[1cm] \\
\hline
\end{tabularx}
\end{center}

\begin{enumerate}[leftmargin=1.5em]
    \item Vérifier que l'effectif total est bien $N = 25$.
    \item Compléter la ligne des fréquences en pourcentage.
    \item Calculer le nombre moyen de couverts servis par jour : $\bar{x} = \frac{\dots \dots}{25} = \pointille[2cm]$.
    \item Déterminer le nombre médian de couverts $Me$.
\end{enumerate}
\lignesreponse[3]
\end{exercice}

\begin{exercice}[3 pts]{\capSApproprier}{Lecture de Diagramme en Bâtons (Garage Automobile)}
Un garage automobile a relevé le nombre d'interventions de vidange effectuées par jour sur une période de 20 jours :

\begin{center}
\begin{tikzpicture}[scale=0.75]
    \draw[step=1cm, slate200, very thin] (0,0) grid (6,7);
    \draw[thick, ->, >=stealth, slate800] (0,0) -- (6.4,0) node[right, font=\scriptsize] {Nombre de vidanges ($x_i$)};
    \draw[thick, ->, >=stealth, slate800] (0,0) -- (0,7.3) node[above, font=\scriptsize] {Nombre de jours ($n_i$)};
    \foreach \x in {1,2,3,4,5}
        \draw (\x,0.1) -- (\x,-0.1) node[below, font=\scriptsize] {\x};
    \foreach \y in {1,2,3,4,5,6,7}
        \draw (0.1,\y) -- (-0.1,\y) node[left, font=\scriptsize] {\y};
    
    % Bâtons
    \draw[line width=8pt, colPrimary] (1,0) -- (1,3);
    \draw[line width=8pt, colPrimary] (2,0) -- (2,6);
    \draw[line width=8pt, colPrimary] (3,0) -- (3,7);
    \draw[line width=8pt, colPrimary] (4,0) -- (4,3);
    \draw[line width=8pt, colPrimary] (5,0) -- (5,1);
\end{tikzpicture}
\end{center}

\begin{enumerate}[leftmargin=1.5em]
    \item Combien de jours le garage a-t-il réalisé exactement 3 vidanges ? \pointille[3cm]
    \item Calculer le nombre total de vidanges réalisées sur les 20 jours.
    \item Calculer le nombre moyen de vidanges par jour.
    \item Quel est le pourcentage de jours où le garage a réalisé au moins 4 vidanges ?
\end{enumerate}
\lignesreponse[3]
\end{exercice}

\begin{exercice}[4 pts]{\capAnalyser}{Comparaison de deux fournisseurs (Logistique \& BTP)}
Une entreprise de travaux publics teste deux marques de sacs de ciment de $25\text{ kg}$ pour vérifier la régularité du remplissage. Elle pèse un échantillon de 10 sacs pour chaque fournisseur :
\begin{itemize}[leftmargin=1.5em, itemsep=1pt]
    \item \textbf{Fournisseur A (en kg)} : $24{,}8$ ; $25{,}0$ ; $25{,}1$ ; $24{,}9$ ; $25{,}2$ ; $25{,}0$ ; $24{,}9$ ; $25{,}1$ ; $25{,}0$ ; $25{,}0$
    \item \textbf{Fournisseur B (en kg)} : $24{,}2$ ; $25{,}8$ ; $24{,}5$ ; $25{,}5$ ; $25{,}0$ ; $24{,}4$ ; $25{,}6$ ; $25{,}0$ ; $24{,}3$ ; $25{,}7$
\end{itemize}

\begin{enumerate}[leftmargin=1.5em]
    \item Calculer la masse moyenne pour chaque fournisseur : $\bar{x}_A = \pointille[2cm]$ et $\bar{x}_B = \pointille[2cm]$.
    \item Calculer l'étendue de chaque série : $e_A = \pointille[2cm]$ et $e_B = \pointille[2cm]$.
    \item Quel fournisseur l'entreprise doit-elle privilégier pour garantir la régularité de ses dosages ? Justifier.
\end{enumerate}
\lignesreponse[4]
\end{exercice}

% ------------------------------------------------------------------------------
% 7. TP TICE / CALCULATRICE NUMWORKS
% ------------------------------------------------------------------------------
\section{TP TICE : Statistiques sur NumWorks}

\begin{tpinfo}[Calculatrice NumWorks]{Application Statistiques}
\begin{enumerate}[leftmargin=1.5em, itemsep=2pt]
    \item \textbf{Saisir la série de données (Application Statistiques)} :
    \begin{itemize}
        \item Ouvrir l'application \textbf{Statistiques} (icône diagramme à barres).
        \item Dans la colonne \textbf{Valeurs V1}, entrer les valeurs : $10$, $12$, $14$, $16$, $18$.
        \item Dans la colonne \textbf{Effectifs N1}, entrer les effectifs : $3$, $8$, $12$, $5$, $2$.
    \end{itemize}
    \item \textbf{Lire les indicateurs statistiques (Onglet Stats)} :
    \begin{itemize}
        \item Aller dans l'onglet \textbf{Stats} en haut de l'écran.
        \item Noter les indicateurs calculés instantanément par la calculatrice :
        \begin{itemize}
            \item Effectif total $n = \pointille[2cm]$
            \item Moyenne $\bar{x} = \pointille[2cm]$
            \item Médiane $Med = \pointille[2cm]$
            \item Étendue = $x_{\text{max}} - x_{\text{min}} = \pointille[2cm]$
        \end{itemize}
    \end{itemize}
    \item \textbf{Visualiser le graphique (Onglet Graphique)} :
    \begin{itemize}
        \item Aller dans l'onglet \textbf{Graphique} $\rightarrow$ choisir \textbf{Histogramme} ou \textbf{Boîte à moustaches}.
    \end{itemize}
\end{enumerate}
\end{tpinfo}

% ------------------------------------------------------------------------------
% 8. VERS LE BREVET (DNB PRO)
% ------------------------------------------------------------------------------
\section{Vers le Brevet (DNB Pro)}

\begin{sujetdnb}[DNB Pro Métropole][10 pts]{Enquête sur les Déplacements des Salariés}
Une grande surface commerciale réalise une enquête auprès de ses \textbf{80 employés} concernant la distance (en kilomètres) séparant leur domicile de leur lieu de travail :

\begin{center}
\renewcommand{\arraystretch}{1.2}
\begin{tabularx}{\linewidth}{|l|*{5}{>{\centering\arraybackslash}X|}}
\hline
\rowcolor{slate800} \textcolor{white}{\textbf{Distance en km ($x_i$)}} & \textcolor{white}{5 km} & \textcolor{white}{10 km} & \textcolor{white}{15 km} & \textcolor{white}{20 km} & \textcolor{white}{30 km} \\
\hline
\textbf{Effectif ($n_i$)} & 24 & 28 & 16 & 8 & 4 \\
\hline
\textbf{Fréquence en \%} & \pointille[1.5cm] & \pointille[1.5cm] & \pointille[1.5cm] & \pointille[1.5cm] & \pointille[1.5cm] \\
\hline
\end{tabularx}
\end{center}

\begin{enumerate}[leftmargin=1.5em, itemsep=2pt]
    \item \capSApproprier\ Vérifier que le nombre total de salariés interrogés est bien égal à $80$.
    \item \capRealiser\ Compléter le tableau en calculant la fréquence en pourcentage pour chaque distance.
    \item \capRealiser\ Calculer la distance moyenne parcourue par les employés :
    \[ \bar{x} = \frac{24 \times 5 + 28 \times 10 + 16 \times 15 + 8 \times 20 + 4 \times 30}{80} = \pointille[3cm] \]
    \item \capAnalyser\ Déterminer la distance médiane de cette série. Interpréter ce résultat.
    \item \capValider\ Calculer le pourcentage de salariés qui habitent à \textbf{moins de 15 km} de leur lieu de travail.
    \item \capCommuniquer\ La direction décide de verser une prime transport aux employés habitant à \textbf{20 km ou plus}. Combien de salariés bénéficieront de cette prime ?
\end{enumerate}

\vspace{1mm}
\noindent\textbf{Rédigez votre démarche ci-dessous :}
\lignesreponse[6]
\end{sujetdnb}
